\listfiles
\documentclass[\OPTclass]{book}
\usepackage[\OPTgeometry]{geometry}
\usepackage{amsmath,amsthm,amssymb}
\usepackage[utf8]{inputenc}
\usepackage[T1]{fontenc}
\usepackage[french]{babel}
\frenchbsetup{
  og=«, fg=»,
  ThinColonSpace=true
}
\addto\captionsfrench{%
  \renewcommand{\contentsname}{Sommaire}%
}

\usepackage{lmodern}
\usepackage[upright,widespace]{fourier}
\usepackage{maths}
\usepackage{microtype}
\usepackage[shortlabels]{enumitem}
\usepackage{booktabs}
\usepackage{pgffor}
\usepackage{graphicx}
\usepackage[x11names]{xcolor}
\usepackage{titlesec}
\usepackage{titletoc}
\usepackage{fancyhdr}

\usepackage[
  colorlinks=true,
  linkcolor=Blue4,
  urlcolor=Blue4,
  linktocpage=true,
  pdftitle={Éléments de théorie des groupes. Solutions des exercices.},
  pdfauthor={Éric Guirbal},
  pdfsubject={Mathématiques},
  pdfkeywords={théorie des groupes, groupe monogène, groupe cyclique, groupe
    diédral, groupe fini, theorèmes de Sylow, sous-groupe normal, action de
    groupe, suite de composition, groupe abélien, groupe libre},
  pdflang=fr-FR]{hyperref}

%%% En-têtes
\pagestyle{fancy}
\fancyhf{}
\renewcommand{\chaptermark}[1]{%
  \markboth{\sffamily#1}{}}
\fancyhead[RE]{\sffamily Éléments de théorie des groupes. Solutions des exercices}
\fancyhead[LO]{\sffamily Chapitre \thechapter. \leftmark}
\fancyhead[RO,LE]{\thepage}
\renewcommand{\headrulewidth}{0pt}

%%% Chapitres
\titleformat{\chapter}[display]
  {\filcenter}
  {\textls[150]{\MakeUppercase{Chapitre\,\, \thechapter}}}
  {0.5pc}  
  {\LARGE\itshape}
  []
\titlespacing{\chapter}
  {0pt}
  {0ex}
  {10ex}

\renewcommand{\thechapter}{%                                                    
  \ifnum\value{chapter}=1
    premier%
  \else
    \Roman{chapter}%
  \fi
}

%%% Sommaire
\titlecontents{chapter}
  [0pc] % left
  {\addvspace{3pt}} % above
  {{\scshape Chapitre \thecontentslabel} / \itshape} % before with label
  {\scshape \thecontentslabel} % before without label
  {\titlerule*[0.5pc]{.}\contentspage} % filler and page
  [] % after

%%% Listes
\setlist[enumerate]{label=\alph*),leftmargin=*,widest=a}

%%% Environnements
\theoremstyle{definition}
\newtheorem{solution}{}
\theoremstyle{remark}
\newtheorem*{remarque}{Remarque}

%%% Pour chacun des chapitres, le nombre d'exercices
%%% et le dossier dans lequel sont leurs solutions.
\newcommand\thebook{%
  % dossier/chapitre/nombre d'exercices
  1-structure_de_groupe/Structure de groupe/33,
  2-classes_modulo_un_sous_groupe/Classes modulo un sous-groupe/10,
  3-groupes_monogenes/{Groupes monogènes, symétriques et diédraux}/33,
  4-sous-groupes_normaux/Sous-groupes normaux/38,
  5-groupe_operant_sur_un_ensemble/Groupe opérant sur un ensemble/27,
  6-groupes_finis/Groupes finis. Théorèmes de Sylow/26,
  7-suites_de_composition/Suites de composition/35,
  8-groupes_abeliens/Groupes abéliens/23,
  9-groupes_libres/Groupes libres. Générateurs et relations. Produit libre de groupes/12
}

\begin{document}

\title{
  Éléments de théorie des groupes\\
  Solutions des exercices
}

\author{Éric Guirbal}
\date{}

\frontmatter

\maketitle

%%%%%%%%%%%%%%%%%%% COPYRIGHT %%%%%%%%%%%%%%%%%%%%%%%%%%%%%%%%%%%%%%%%%%%%%%%%
\input{revision}
\thispagestyle{empty}
\null
\vfill
\begin{center}
  \small

  \textbf{%
    Éléments de théorie des groupes \\
    Solutions des exercices
  }
  \medskip

  Éric \textsc{Guirbal}

  \medskip

  Version: \gitDescribe{} (\gitDate)

  Compilé le \today

  \vfill

  \includegraphics[width=2cm]{images/cc-by-nc-sa} \\ [0.15cm]

  Ce document est distribué selon les termes de la licence Creative Commons
  Attribution - Pas d'utilisation commerciale - Partage à l'identique 3.0
  France. \\
  \url{https://creativecommons.org/licenses/by-nc-sa/3.0/fr/}
\end{center}


\tableofcontents

\mainmatter

\foreach \dir / \chaptertitle / \numberofexercises [count=\chapternum] in \thebook {
  \chapter{\chaptertitle}
  \foreach \sol in {1,...,\numberofexercises} {%
    \edef\solutionfilename{
      \ifnum \sol < 10
        book/\dir/\chapternum-0\sol.tex
      \else
        book/\dir/\chapternum-\sol.tex
      \fi
    }
    \IfFileExists{\solutionfilename}{%
      \setcounter{solution}{\numexpr\sol-1\relax}
      \begin{solution}
        \input{\solutionfilename}
      \end{solution}
    }{}
  }
}

\backmatter

\bibliographystyle{alpha}
\bibliography{book/bibliographie}
\nocite{*}
\end{document}

